% Part: lambda-calculus
% Chapter: lambda-definability
% Section: arithmetical-functions

\documentclass[../../../include/open-logic-section]{subfiles}

\begin{document}

\olfileid{lam}{rep}{arf}
\olsection{الدوال الحسابية \usetoken{S}{lambda definable}}

\begin{prop}
  \ollabel{prop:succ-ld}
  دالة الخلف~$\Succ$ !!{lambda definable}.
\end{prop}

\begin{proof}
الحد الذي !!{lambda define}s دالة الخلف هو
\begin{align*}
  \fn{Succ} & \ident \lambd[a][\lambd[fx][f ({a} f x)]].
  \intertext{وهذا، وفق اصطلاحاتنا، اختصار لـ}
  \fn{Succ} & \ident \lambd[a][\lambd[f][\lambd[x][(f (({a} f) x))]]].
  \intertext{إن~$\fn{Succ}$ دالة تقبل عددًا~$a$ وسيطًا، وتقيّم إلى
    دالة أخرى هي~$\lambd[fx][f ({a} f x)]$. وهذه الدالة ليست هي نفسها
    عدد تشيرش. غير أنها تختزل إلى عدد تشيرش إذا كان الوسيط~$a$ عدد
    تشيرش. انظر:}
   (\lambd[a][\lambd[fx][f ({a} f x)]])\,\num{n} & \redone
   \lambd[fx][f ({\num{n}} f x)].
   \intertext{الحد المضمّن~$\num{n}fx$ موضع اختزال، لأن~$\num{n}$ هو
     $\lambd[fx][f^nx]$. ولذلك~$\num{n}fx \red f^nx$، ومن ثم لدينا
     للحد كله}
   \fn{Succ}\, \num n & \red \lambd[fx][f(f^n(x))],
\end{align*}
أي~$\num{n+1}$.
\end{proof}

\begin{ex}
  لننظر فيما يحدث عندما نطبق~$\fn{Succ}$ في~$\num{0}$، أي في
  $\lambd[fx][x]$. وسنكتب الحدود كاملة:
  \begin{align*}
    \fn{Succ}\, \num{0} & \ident (\lambd[a][\lambd[f][\lambd[x][(f (({a} f) x))]]])(\lambd[f][\lambd[x][x]])\\
    & \redone \lambd[f][\lambd[x][(f (({(\lambd[f][\lambd[x][x]])} f) x))]] \\
    & \redone \lambd[f][\lambd[x][(f ({(\lambd[x][x])} x))]] \\
    & \redone \lambd[f][\lambd[x][(f x)]] \ident \num{1}\\
  \end{align*}
\end{ex}

\begin{prob}
  الحد
  \begin{align*}
    \fn{Succ}' & \ident \lambd[{n}][\lambd[fx][{n} f (f x)]]
  \end{align*}
  !!{lambda define}s دالة الخلف. اشرح السبب.
\end{prob}

\begin{prop}
  \ollabel{prop:add-ld}
  دالة الجمع~$\Add$ !!{lambda definable}.
\end{prop}

\begin{proof}
يمكن أن !!{lambda define} الجمع بالحدين
\begin{align*}
  \fn{Add} & \ident \lambd[{a}{b}][\lambd[fx][{a} f ({b} f x)]]
\intertext{أو، بدلًا من ذلك،}
  \fn{Add}' & \ident \lambd[{a}{b}][{a}\, \fn{Succ} \, {b}].
\intertext{يعمل التمثيل الأول للجمع كما يأتي: تقبل~$\fn{Add}$ أولًا
  عددين~${a}$ و${b}$. والناتج دالة تقبل~$f$ و$x$ وتعيد~$af(bfx)$.
  فإذا كان~$a$ و$b$ عددي تشيرش~$\num{n}$ و$\num{m}$، اختزل ذلك إلى
  $f^{n+m}(x)$، وهو مطابق لـ$f^{n}(f^{m}(x))$. وبالتفصيل:}
  (\lambd[{a}{b}][\lambd[fx][{a} f ({b} f x)]])\num n\,\num m & \red
  \lambd[fx][\num{n}\, f (\num {m}\, f x)] \\
  & \red \lambd[fx][\num{n}\, f (f^m x)] \\
  & \red \lambd[fx][f^n (f^m x)] \ident \num {n+m}.
\intertext{أما التمثيل الثاني للجمع~$\fn{Add'}$ فيعمل على نحو مختلف.
  فعند تطبيقه في عددي تشيرش~$\num{n}$ و$\num{m}$،}
\fn{Add}' \num n \,\num m
& \red \num{n}\, \fn{Succ}\, \num{m}.
\intertext{لكن~$\num{n} f x$ يختزل دائمًا إلى~$f^n(x)$. لذلك،}
  \num{n}\,  \fn{Succ}\, \num{m} & \red \fn{Succ}^n(\num{m}).
\end{align*}
وبما أن~$\fn{Succ}$ !!{lambda define}s دالة الخلف، وأن تطبيق دالة
الخلف~$n$ مرة في~$m$ يعطي~$n+m$، فإن هذا بدوره يختزل إلى~$\num{n+m}$.
\end{proof}

\begin{prop}
  \ollabel{prop:mult-ld}
  الضرب !!{lambda definable} بالحد
  \[
  \fn{Mult} \ident \lambd[ab][\lambd[fx][a (b f) x]]
  \]
\end{prop}

\begin{proof}
  لبيان كيفية عمله، افترض أننا نطبق~$\fn{Mult}$ في عددي تشيرش
  $\num{n}$ و$\num{m}$. يختزل~$\fn{Mult} \, \num{n} \, \num{m}$
  إلى~$\lambd[fx][\num{n}(\num{m}\, f)x]$. ويعرّف الحد~$\num{m} f$
  دالة تطبق~$f$ في وسيطها~$m$ مرة. ومن ثم فإن
  $\num{n} (\num{m} f) x$ يطبق الدالة ``طبق~$f$ عدد~$m$ من المرات''
  نفسها~$n$ مرة في~$x$. وبعبارة أخرى، نطبق~$f$ في~$x$ عدد~$n\cdot m$
  من المرات. لكن الحد السوي الناتج هو تحديدًا عدد تشيرش~$\num{nm}$.
\end{proof}

\begin{editorial}
يمكننا في الواقع تبسيط هذا الحد أكثر باختزال~$\eta$:
\[
  \fn{Mult} \ident \lambd[ab][\lambd[f][a (b f)]].
\]

لكن علينا عندئذ أن نشرح اختزال~$\eta$ أولًا.
\end{editorial}

\begin{prob}
يمكن أن !!{lambda define} الضرب بالحد
\[
  \fn{Mult}' \ident \lambd[ab][a (\fn{Add}\, b) \num{0}].
\]
اشرح لماذا يعمل هذا الحد.
\end{prob}

إن تعريف الأس بحد~$\lambd$ بسيط على نحو مدهش:
\[
  \fn{Exp} \ident \lambd[be][e b].
\]
الوسيط الأول~$b$ هو الأساس، والثاني~$e$ هو الأس. وحدسيًا، يكون~$e f$
هو~$f^e$ بحسب ترميزنا للأعداد. وإذا استعصى فهم ذلك، فلا يزال بوسعنا
تعريف الأس بالضرب المكرر أيضًا:
\[
  \fn{Exp}' \ident \lambd[be][e (\fn{Mult}\, b) \num{1}].
\]

ليست دالة السابق والطرح على أعداد تشيرش ببساطة ما قد نظنه؛ إذ يتطلبان
ترميز الأزواج.

\end{document}
